% !TeX spellcheck = en_GB
\documentclass[a4paper,12pt,titlepage]{report}
\usepackage[
    twoside=true,
    left=3cm, % inner
    right=2cm, % outer
    top=2.5cm,
    bottom=2.5cm,
]{geometry}
\linespread{1.5}
\usepackage[skip=4pt]{parskip}

\usepackage[utf8]{inputenc}
\usepackage[T1]{fontenc}
\usepackage[english,polish]{babel}
\addto\captionspolish{\renewcommand{\figurename}{Rys.}}

\usepackage{graphicx}
\usepackage{amsfonts}
\usepackage{amsmath}
\usepackage{amssymb}
\usepackage{wasysym}
\usepackage{scalerel}
\usepackage{siunitx}
\usepackage{textcomp}
\usepackage{gensymb}
\usepackage{parskip}
\usepackage{titlesec}
\usepackage{listings}
\usepackage{lipsum}
\usepackage[table,xcdraw]{xcolor}
\usepackage{enumitem}
\usepackage[font=small,labelfont=bf]{caption}
\usepackage{hyperref}
\usepackage{csquotes}
\usepackage{textgreek}
\usepackage{url}
\usepackage{rotating}
\usepackage{fancyhdr}
\usepackage{lastpage}
\usepackage{bm}
\usepackage{multirow}
\usepackage{subcaption}
\usepackage[export]{adjustbox}

\usepackage{tikz}
\usetikzlibrary{positioning}
\usetikzlibrary{external}
\tikzexternalize[prefix=tikz/]

% Highlight thesis author in bibliobraphy
% \usepackage{xstring}
% \let\originalbibitem\bibitem
% \def\bibitem#1#2\par{%
%   \noexpandarg
%   \originalbibitem{#1}
%   \StrSubstitute{#2}{I.Nazwisko}{\textbf{I.Nazwisko}}\par}

% ]


%%%%%%%%%%%%%%%%%%%%%%%%%%%%%%%%%%%%%%%%

\linespread{1.5}

\titleformat{\chapter}{\LARGE\bf}{\thechapter}{25pt}{\LARGE\bf}
\titlespacing*{\chapter}{0pt}{-50pt}{10pt}

\titleformat*{\section}{\Large\bf}
\titleformat*{\subsection}{\large\bf}

\renewcommand*\rmdefault{ptm}
\renewcommand*\sfdefault{phv}
\renewcommand*\ttdefault{pcr}

\definecolor{mygray}{rgb}{0.97,0.97,0.97}
\lstset{basicstyle=\small\ttfamily, frame=single, backgroundcolor=\color{mygray}}
\urlstyle{same}

\widowpenalty 1800
\clubpenalty 1800


\setcounter{tocdepth}{2}
% \setcounter{secnumdepth}{2}

\graphicspath{ {} }

\setlength{\headheight}{14.49998pt}
\addtolength{\topmargin}{-2.49998pt}

\fancyhf{}
\rhead{Szablon pisania pracy dyplomowej}
\lhead{inż. Imię Nazwisko}
\fancyfoot[RO,LE]{\thepage} % strona tytułowa drukowana jednostronnie
% \fancyfoot[RE,LO]{\thepage} % strona tytułowa drukowana dwustronnie

\fancypagestyle{plain}{} % to add header in special pages


\begin{document}
%%%%%%%%%%%%%%%%%%%%%%%%%%%%%%%%%%%%%%%%

\begin{titlepage}
    \includegraphics[height=3cm, right]{images/LUT_wydzial.png}
    \hfill
    
	\centering
    \vspace{24pt}
    
	{
        \large 
        inż. Imię Nazwisko \\
        123456 % numer albumu
        \par
    }

	\vfill

	{\Large PRACA DYPLOMOWA MAGISTERSKA \\
    \large na kierunku Sztuczna Inteligencja i Uczenie Maszynowe}\par
	% \vspace{24pt}
    \vfill
	{
		\Large\bfseries
        Szablon pisania pracy dyplomowej
	    \par
	}

	\vfill


	\raggedright
    {
        \large
        Instytut Informatyki Stosowanej
        \par
    }
	{
        \large 
        Promotor: 
		dr hab. inż. Imię Nazwisko
        \par
    }
    {
        \large 
        Promotor pomocniczy: 
		dr inż. Imię Nazwisko 
        \par
    }
    \vspace{24pt}

	\centering
	{\large Łódź 2024 \par}

\end{titlepage}

%%%%%%%%%%%%%%%%%%%%%%%%%%%%%%%%%%%%%%%%
\pagestyle{fancy}
\pagenumbering{roman}
% \setcounter{page}{2}

\newenvironment{abstractpage}
{\cleardoublepage}
{\vfill\cleardoublepage}

\renewenvironment{abstract}[1]
{\smallskip\selectlanguage{#1}%
	\begin{center}\bfseries\abstractname\end{center}}
{\par\smallskip}

\begin{abstractpage}
	\begin{abstract}{polish}
        Lorem ipsum dolor sit amet, consectetur adipiscing elit, sed do eiusmod tempor incididunt ut labore et dolore magna aliqua. Odio pellentesque diam volutpat commodo sed. Proin fermentum leo vel orci porta non pulvinar neque laoreet. Diam volutpat commodo sed egestas egestas fringilla. Elit at imperdiet dui accumsan sit amet nulla. Vel pretium lectus quam id leo in vitae. Faucibus vitae aliquet nec ullamcorper sit amet risus. Leo vel fringilla est ullamcorper eget nulla facilisi etiam dignissim. Quis risus sed vulputate odio ut enim blandit volutpat. Tempor commodo ullamcorper a lacus vestibulum. Tincidunt augue interdum velit euismod. Ullamcorper dignissim cras tincidunt lobortis feugiat. A condimentum vitae sapien pellentesque habitant morbi tristique senectus. Volutpat ac tincidunt vitae semper quis lectus nulla at. Tortor aliquam nulla facilisi cras fermentum odio eu feugiat. Leo integer malesuada nunc vel risus. Turpis in eu mi bibendum neque egestas congue quisque. Mattis rhoncus urna neque viverra justo nec ultrices dui. Lorem sed risus ultricies tristique nulla. Vitae tortor condimentum lacinia quis vel eros.
		
		Słowa kluczowe: 
	\end{abstract}
	\pagebreak
	\begin{abstract}{english}
         Lorem ipsum dolor sit amet, consectetur adipiscing elit, sed do eiusmod tempor incididunt ut labore et dolore magna aliqua. Odio pellentesque diam volutpat commodo sed. Proin fermentum leo vel orci porta non pulvinar neque laoreet. Diam volutpat commodo sed egestas egestas fringilla. Elit at imperdiet dui accumsan sit amet nulla. Vel pretium lectus quam id leo in vitae. Faucibus vitae aliquet nec ullamcorper sit amet risus. Leo vel fringilla est ullamcorper eget nulla facilisi etiam dignissim. Quis risus sed vulputate odio ut enim blandit volutpat. Tempor commodo ullamcorper a lacus vestibulum. Tincidunt augue interdum velit euismod. Ullamcorper dignissim cras tincidunt lobortis feugiat. A condimentum vitae sapien pellentesque habitant morbi tristique senectus. Volutpat ac tincidunt vitae semper quis lectus nulla at. Tortor aliquam nulla facilisi cras fermentum odio eu feugiat. Leo integer malesuada nunc vel risus. Turpis in eu mi bibendum neque egestas congue quisque. Mattis rhoncus urna neque viverra justo nec ultrices dui. Lorem sed risus ultricies tristique nulla. Vitae tortor condimentum lacinia quis vel eros.

		Keywords: 
	\end{abstract}
\end{abstractpage}

%%%%%%%%%%%%%%%%%%%%%%%%%%%%%%%%%%%%%%%%
\pagestyle{fancy}
% \pagenumbering{roman}
\tableofcontents

%%%%%%%%%%%%%%%%%%%%%%%%%%%%%%%%%%%%%%%%
\chapter{Wprowadzenie} \label{ch:Intro}
\pagenumbering{arabic}

\section{Przykładowe obrazki}
\begin{figure}[h]
    \centering
    \includegraphics[width=0.75\textwidth]{images/Introduction/PittsMcCulloch.png}
    \caption{Przykładowe umieszczenie pojedynczego obrazu}
    \label{fig:Intro:PittsMcCulloch}
\end{figure}

\begin{figure}[h]
     \centering
     \begin{subfigure}[b]{0.45\linewidth}
         \centering
         \includegraphics[width=\linewidth]{images/Introduction/PittsMcCulloch.png}
         \caption{
         Lewy obrazek
         }
         \label{fig:Intro:PittsMcCulloch:lewy}
     \end{subfigure}
     \hfill
     \begin{subfigure}[b]{0.45\linewidth}
         \centering
         \includegraphics[width=\linewidth]{images/Introduction/PittsMcCulloch.png}
         \caption{
         Prawy obrazek
         }
         \label{fig:Intro:PittsMcCulloch:prawy}
     \end{subfigure}
        \caption[Przykładowe umieszczenie kilku obrazów]{
        Dłuższy opis do przykładowego umieszczenia kilku obrazów.
        }
\end{figure}

\begin{figure}[h]
    \centering
    \begin{tikzpicture}[scale=\textwidth/16cm, samples=200]
    % \draw[step=1cm,gray,very thin] (0, -13) grid (13, 0); % Linie pomocnicze do wyznaczania pozycji
    
    \node[align=center] (image) at (0, -2) {Image};
    
    \draw[->] (image.east) -- (1, -2);
    
    \draw (1, 0) -- (3, -1) -- (3, -3) -- (1, -4) -- cycle; % Encoder
    \draw (2, -2) node [align=center] {Encoder};
    
    \draw[->] (3, -2) -- (3.5, -2);
    
    \begin{scope}[color=teal] % Latent
        \draw (3.5, -1) -- (3.5, -3) -- (4, -3) -- (4, -1) -- cycle;
        \draw (3.75, -0.75) node [align=center] {Latent};
        \draw (3.75, -3.25) node [align=center] {$\in \mathbb{R}^{1 \times w}$};
    \end{scope}

    \draw[->] (4, -2) -- (4.5, -2);

    \draw (4.5,-1) -- (6.5, 0) -- (6.5, -4) -- (4.5, -3) -- cycle; % Decoder
    \draw (5.5, -2) node [align=center] {Decoder};

    \node[align=center] (image) at (7.5, -2) {Image};

    \draw[->] (6.5, -2) -- (image.west);   

    \end{tikzpicture}
    \caption{Przykładowy schemat w tikz}
    \label{fig:Intro:AutoEncoder}
\end{figure}

\section{Przykładowe równania}
Przykładowe równanie przedstawiono na \eqref{eq:Intro:DeltaRule},
\begin{equation} \label{eq:Intro:DeltaRule}
    \Delta \vec{w_j} = \eta \epsilon_j f'(s_j) \vec{x_j} = \eta (t_j - y_j) f'(\vec{w_{j-1}}^T \vec{x_j}) \vec{x_j}
\end{equation}
podczas gdy zbiór równań przedstawiono na \eqref{eq:Intro:GiniEntropy}
\begin{equation} \label{eq:Intro:GiniEntropy}
    \begin{split}
        H = - \sum_i p_i log(p_i) \\
        G = 1 - \sum_i p_i^2.
    \end{split}
\end{equation}

\section{Przykładowa tabela}
% wygenerowane przy pomocy https://www.tablesgenerator.com/
\begin{table}[h]
\centering
\caption{Przykładowa tabela}
\label{tab:Intro:PrzykladowaTabela}
\begin{tabular}{c|c|c}
         & \textbf{Kolumna 1} & \textbf{Kolumna 2} \\ \hline
Wiersz 1 & 1                  & 2                  \\
Wiersz 2 & 3                  & 4                 
\end{tabular}
\end{table}

%%%%%%%%%%%%%%%%%%%%%%%%%%%%%%%%%%%%%%%%
\chapter{Cel i zakres pracy}
Lorem ipsum dolor sit amet, consectetur adipiscing elit, sed do eiusmod tempor incididunt ut labore et dolore magna aliqua. Odio pellentesque diam volutpat commodo sed. Proin fermentum leo vel orci porta non pulvinar neque laoreet. Diam volutpat commodo sed egestas egestas fringilla. Elit at imperdiet dui accumsan sit amet nulla. Vel pretium lectus quam id leo in vitae. Faucibus vitae aliquet nec ullamcorper sit amet risus. Leo vel fringilla est ullamcorper eget nulla facilisi etiam dignissim. Quis risus sed vulputate odio ut enim blandit volutpat. Tempor commodo ullamcorper a lacus vestibulum. Tincidunt augue interdum velit euismod. Ullamcorper dignissim cras tincidunt lobortis feugiat. A condimentum vitae sapien pellentesque habitant morbi tristique senectus. Volutpat ac tincidunt vitae semper quis lectus nulla at. Tortor aliquam nulla facilisi cras fermentum odio eu feugiat. Leo integer malesuada nunc vel risus. Turpis in eu mi bibendum neque egestas congue quisque. Mattis rhoncus urna neque viverra justo nec ultrices dui. Lorem sed risus ultricies tristique nulla. Vitae tortor condimentum lacinia quis vel eros.


Lorem ipsum dolor sit amet, consectetur adipiscing elit, sed do eiusmod tempor incididunt ut labore et dolore magna aliqua. Odio pellentesque diam volutpat commodo sed. Proin fermentum leo vel orci porta non pulvinar neque laoreet. Diam volutpat commodo sed egestas egestas fringilla. Elit at imperdiet dui accumsan sit amet nulla. Vel pretium lectus quam id leo in vitae. Faucibus vitae aliquet nec ullamcorper sit amet risus. Leo vel fringilla est ullamcorper eget nulla facilisi etiam dignissim. Quis risus sed vulputate odio ut enim blandit volutpat. Tempor commodo ullamcorper a lacus vestibulum. Tincidunt augue interdum velit euismod. Ullamcorper dignissim cras tincidunt lobortis feugiat. A condimentum vitae sapien pellentesque habitant morbi tristique senectus. Volutpat ac tincidunt vitae semper quis lectus nulla at. Tortor aliquam nulla facilisi cras fermentum odio eu feugiat. Leo integer malesuada nunc vel risus. Turpis in eu mi bibendum neque egestas congue quisque. Mattis rhoncus urna neque viverra justo nec ultrices dui. Lorem sed risus ultricies tristique nulla. Vitae tortor condimentum lacinia quis vel eros.

\newpage

%%%%%%%%%%%%%%%%%%%%%%%%%%%%%%%%%%%%%%%%
\chapter{Przegląd literatury}
W pracy dyplomowej należy cytować źródła, zwłaszcza artykuły naukowe \cite{vaswani2017, dosovitskiy2020}.
Warto także zacytować wykorzystywane biblioteki programistyczne \cite{paszke}.

\newpage

%%%%%%%%%%%%%%%%%%%%%%%%%%%%%%%%%%%%%%%%
\chapter{Materiały i metody}
\newpage

%%%%%%%%%%%%%%%%%%%%%%%%%%%%%%%%%%%%%%%%
\chapter{Wyniki}
\newpage

%%%%%%%%%%%%%%%%%%%%%%%%%%%%%%%%%%%%%%%%
\chapter{Podsumowanie}
\newpage

%%%%%%%%%%%%%%%%%%%%%%%%%%%%%%%%%%%%%%%%
% \addcontentsline{toc}{chapter}{Podziękowania}
% \chapter*{Podziękowania}
% Część badań przestawionych w tej pracy została sfinansowana z \dots
% \newpage

%%%%%%%%%%%%%%%%%%%%%%%%%%%%%%%%%%%%%%%%

\newpage
\addcontentsline{toc}{chapter}{\listfigurename}
\listoffigures

\newpage
\addcontentsline{toc}{chapter}{\listtablename}
\listoftables
%%%%%%%%%%%%%%%%%%%%%%%%%%%%%%%%%%%%%%%%

\newpage
\chapter*{Wykaz symboli i oznaczeń}
\addcontentsline{toc}{chapter}{Wykaz symboli i oznaczeń}
O ile nie jest to inaczej wyszczególnione w tekście praca ta wykorzystuje poniższą notacje:
\begin{itemize}
    \item $a$ (mała litera kursywą) oznacza skalar,
    \item $\bm{a}$ (mała pogrubiona litera) oznacza wektor, 
    \item $\bm{A}$ (wielka pogrubiona litera) oznacza macierz,
    \item $\bm{A_i}$ oznacza $i$-ty wiersz macierzy - wektor,
    \item $a_i$ oznacza $i$-ty element wektora - skalar,

    \item $\mathbb{B}$ oznacza zbiór liczb binarnych $\{0, 1\}$,
    \item $\pm \mathbb{B}$ oznacza zbiór liczb bipolarnych $\{-1, 1\}$
    .
\end{itemize}

\newpage
\chapter*{Wykaz używanych skrótów}
\addcontentsline{toc}{chapter}{Wykaz używanych skrótów}
\begin{itemize}
    \item IEEE - Institute of Electrical and Electronics Engineers,
    \item ACM - Association for Computing Machinery.
\end{itemize}

\clearpage
\addcontentsline{toc}{chapter}{Bibliografia}
\bibliographystyle{apalike}
\bibliography{MyLibrary}

%%%%%%%%%%%%%%%%%%%%%%%%%%%%%%%%%%%%%%%%

\end{document}
